\begin{lemma}{Erwartungswertlemma}{Erwartungswertlemma}
Seien $f,g\in L_1(\Omega,\mfa,\Prim,\R)$ mit $f,g\ge 0$ und $\E(f)\le\E(g)$.\\
Seien $a, b \in \R_{>0}$ mit $a \le \E(f) \le \E(g) \le b$.\\
Dann ist $\meng{f > 0}\cap \meng{ag \le bf} \in \mfa$ und es gilt:
$\Prim\big(\meng{f > 0}\cap \meng{ag \le bf}\big) > 0\text.$\\
Insbesondere existiert also ein $\omega \in \Omega$ mit $f(\omega) > 0$ und
$\frac {g(\omega)}{f(\omega)} \le \frac b a$.
\end{lemma}
\begin{proof}
Da $f$ und $g$ integrierbar sind, sind auch $ag$ und $bf$ integrierbar. Damit
sind sie messbar und daher sind $\meng{f >0}$,
$\meng{ag \le bf} \in \mfa$. Also auch deren Schnitt.\\
Annahme: $\Prim(\meng{f > 0}\cap \meng{ag \le bf}) = 0$.\\
Dann ist $\Prim(f = 0 \vee ag > bf) = 1$.\\
Es ist 
\begin{align*}\meng{f = 0 \vee ag > bf} &= \meng{f=0} \cup \meng{ag > bf} \\
&= \meng{f=0} \overset{\cdot}\cup \bigg(\meng{ag>bf}\setminus\meng{f=0}
\bigg)\text.\end{align*}
Da $\E(f) > 0$ und $f \ge 0$ ist $\Prim(f = 0) < 1$.\\
Wegen der Disjunktheit von $\meng{ag >bf}\setminus\meng{f=0}$ und $\meng{f=0}$
gilt daher
\[\Prim\left(ag > bf \wedge f > 0 \right) = 1 - \Prim(\meng{f= 0}) > 0\text.\]
Sei $A \coloneqq  \meng{ag>bf}$.\\
Dann ist $\meng{ag>bf}\setminus\meng{f=0}\subseteq A$ und daher $\Prim(A)>0$.\\
Wegen $b>0$ ist $A = \meng{f < \frac a b g}$.\\
Seien $X \coloneqq  f \Eins_A$ und $Y\coloneqq  \left(\frac a b g\right) \Eins_A$.\\
Dann sind $0 \le X \le f$ und $0 \le Y \le \frac a b g$, also nach
\ref{Integrierbarkeitskriterien}: $X, Y$ integrierbar.\\
Weiter sind $X \le Y$ und $\meng{X < Y} = A$.\\
Nach~\ref{StrengeMonotonieDesErwartungswerts} gilt somit: $\E(X) < \E(Y)$.\\
Erhalte wegen der Integrierbarkeit von $g$:
\begin{align*}
& \E(f) = \E\left(f\Eins_{\meng{f = 0 \vee ag > bf}}\right) 
= \E\left(f\Eins_{\meng{f=0}} +f\Eins_{\meng{ag > bf}\setminus\meng{f=0}}\right)\\
=&\, \E\left(f\Eins_{\meng{ag > bf}\setminus\meng{f=0}} \right) 
\le \E\left(f\Eins_A \right)= \E(X) < \E(Y) = \E\left(\frac abg \right)
=\frac a b \E\left(g\right) \le a
\end{align*}
Was aber wegen $\E(f) \ge a$ nicht sein kann. $\lightning$
\end{proof}